\documentclass[11pt]{article}
\usepackage[margin=26mm]{geometry}
\usepackage{amsmath,amssymb,amsthm}
\usepackage[hidelinks]{hyperref}
\newtheorem{theorem}{Theorem}
\newcommand{\Id}{\operatorname{Id}}
\title{Conjecture 2207: proportion one\protect\\for a finite fraction field}
\author{Local proof package}
\date{8 October 2026}
\begin{document}
\maketitle
\noindent\textbf{Source and interpretation.}
The source says: ``The proportion of 2-generated ideals of a Bezout domain is 1
under finite-quotient-field conditions (2-generated Bezoutization).''
We use the usual meaning of \emph{quotient field}: the field of fractions.
Here \emph{finite} means finitely many field elements. A two-generated ideal
means an ideal of the form $(a,b)$; repeated or redundant generators are allowed.
In particular $(0,0)=(0)$ and $(1,0)=R$ qualify.

\medskip
\noindent For a commutative integral domain $R$, write $\Id(R)$ for its actual
set of ideals and put
\[
 \Id_2(R)=\{I\in\Id(R):\text{there are }a,b\in R\text{ with }I=(a,b)\}.
\]
When $\Id(R)$ is finite, define its two-generated proportion by
$\#\Id_2(R)/\#\Id(R)$. Ideals, rather than presentations by generator pairs,
are counted.

\begin{theorem}
Let $R$ be a B\'ezout integral domain and let $K$ be its field of fractions.
If $K$ has finitely many elements, then $R$ is a field, $\Id(R)$ consists
of exactly the two ideals $(0)$ and $R$, and
\[
 \#\Id_2(R)=\#\Id(R)=2,\qquad
 \frac{\#\Id_2(R)}{\#\Id(R)}=1.
\]
\end{theorem}

\begin{proof}
The fraction-field map $R\longrightarrow K$ is injective. Consequently $R$
is finite; this is a consequence of finite $K$, not an additional hypothesis.
For $a\in R\setminus\{0\}$, the map $x\mapsto ax$ is injective because $R$
is an integral domain. An injective self-map of a finite set is surjective,
so some $b\in R$ satisfies $ab=1$. Thus $R$ is a field.

Let $I$ be an ideal. If $I$ contains a nonzero element $a$, it also contains
$a^{-1}a=1$, and hence $I=R$. Otherwise $I=(0)$. The ideals $(0)$ and $R$
are distinct because a domain has $0\ne1$. They are generated by $(0,0)$
and $(1,0)$, respectively. Every ideal is therefore two-generated, the
numerator and denominator are both $2$, and their quotient is $1$.
\end{proof}

\newpage
\section*{Formal correspondence and scope}
The Lean file uses the existing mathlib objects \texttt{Ideal R},
\texttt{IsBezout R}, and \texttt{IsFractionRing R K}.
The original theorem has precisely the domain, B\'ezout, fraction-field and
finite-field hypotheses:
\begin{center}
\small
\begin{tabular}{ll}
Ring/domain & \texttt{[CommRing R] [IsDomain R]}\\
B\'ezout property & \texttt{[IsBezout R]}\\
Fraction field & \texttt{[Field K] [Algebra R K]}\\
 & \texttt{[IsFractionRing R K]}\\
Finite cardinality of $K$ & \texttt{[Finite K]}.
\end{tabular}
\end{center}
No \texttt{Finite R} or arbitrary \texttt{Fintype} object is a premise of
the original conclusion. Finiteness of $R$ and of its ideal type is proved
from these hypotheses. The real mathlib B\'ezout class asserts that finitely
generated ideals are principal; no replacement axiom is introduced. Its
assumption is retained in the final theorem, although the proof establishes
the stronger result for every domain with finite fraction field.

\medskip
\noindent\textbf{Definitions.}
\texttt{TwoGenerated R I} is the proposition
\[
 \exists a,b\in R,\quad I=\operatorname{span}\{a,b\}.
\]
The numerator counts the filter of all actual ideals by this predicate.
The denominator counts all actual ideals. The ratio is rational-valued:
\begin{center}
\small
\begin{tabular}{@{}ll@{}}
Numerator & \texttt{twoGeneratedCount R}\\
Denominator & \texttt{Nat.card (Ideal R)}\\
Ratio & \texttt{twoGeneratedProportion R}.
\end{tabular}
\end{center}
An enumeration of ideals is obtained only after proving their finiteness.

\medskip
\noindent\textbf{Proof components.}
\begin{center}
\small
\begin{tabular}{@{}p{0.61\linewidth}p{0.34\linewidth}@{}}
\texttt{finite\_domain\_of\_finite\_fraction\_field} & Derives finite $R$.\\
\texttt{isField\_of\_finite\_fraction\_field} & Constructs inverses.\\
\texttt{ideal\_eq\_bot\_or\_top} & Classifies all ideals.\\
\texttt{all\_ideals\_two\_generated} & Gives actual pairs.\\
\texttt{ideal\_card\_eq\_two} & Denominator is $2$.\\
\texttt{twoGeneratedCount\_eq\_two} & Numerator is $2$.\\
\texttt{original\_proportion\_one} & Proportion is $1$.
\end{tabular}
\end{center}
A concrete $R=K=\mathbb F_2$ fraction-field example is also checked.

\medskip
\noindent\textbf{Interpretation limits.}
Finite extension degree over a specified base field and finite residue fields
are different conditions; neither is supplied in the source or used here.
Requiring \emph{minimal} generator number exactly two is also a different
convention from being generated by a pair. No asymptotic density or ideal
ordering is introduced: the fraction-field hypothesis makes the full set of
ideals finite, and the stronger universal generator result applies to it.

\medskip
\noindent\textbf{Provenance.}
The complete fraction-field bridge was first checked in this run's scout
task and reused in the author package. The source is retained byte-for-byte
as \texttt{00000002207.md}. Its SHA256 is
\begin{center}
\small\texttt{0016cb2bf81d535588320e63545be9089d320a430edb18bbe792983b12efb23a}.
\end{center}
This document records a local proof package; acceptance and any submission
eligibility require the separate manager-owned checks and independent review.
\end{document}
